\documentclass{article}

%-------------------------%
% PREAMBUŁA --------------%
%-------------------------%

\usepackage[T1]{fontenc}
\usepackage[utf8]{inputenc}
\usepackage{amsmath}
%\usepackage{amsthm}
%\theoremstyle{plain}
%\usepackage{enumitem}
\usepackage[OT4]{polski}
%\usepackage{caption}
%\usepackage{subcaption}
\usepackage{tabularx}
\usepackage{caption}
%\captionsetup[table]{skip=-4pt}
%\usepackage{tabulary}
%\usepackage{array}
\usepackage{hyperref}
%\usepackage[pdftex]{graphicx}
%\usepackage{sidecap}
%\usepackage{gensymb}

\usepackage{wrapfig}

\newlength{\RoundedBoxWidth}
\newsavebox{\GrayRoundedBox}
\newenvironment{GrayBox}[1][\dimexpr\textwidth-4.5ex]%
   {\setlength{\RoundedBoxWidth}{\dimexpr#1}
    \begin{lrbox}{\GrayRoundedBox}
       \begin{minipage}{\RoundedBoxWidth}}%
   {   \end{minipage}
    \end{lrbox}
    \begin{center}
    \begin{tikzpicture}%
       \draw node[draw=black,fill=black!1,rounded corners,%
             inner sep=2ex,text width=\RoundedBoxWidth]%
             {\usebox{\GrayRoundedBox}};
    \end{tikzpicture}
    \end{center}}

%------- Typy kolumn do ustawiania szerokości ------- %
\newcolumntype{L}[1]{>{\raggedright\let\newline\\\arraybackslash\hspace{0pt}}m{#1}}
\newcolumntype{C}[1]{>{\centering\let\newline\\\arraybackslash\hspace{0pt}}m{#1}}
\newcolumntype{R}[1]{>{\raggedleft\let\newline\\\arraybackslash\hspace{0pt}}m{#1}}

\usepackage[justification=centering]{caption}
\usepackage{graphicx}
\usepackage{lastpage}
\usepackage{fancyhdr}
\usepackage{gensymb}
\usepackage{tikz}
\usepackage{lipsum}
%\pagestyle{fancy}

\usepackage[margin=60pt]{geometry}



\begin{document}

%-------------------------%
%Tabela nagłówkowa -------%
%-------------------------%

\begin{table}[h]
\begin{tabular}{|l|l|l|l|l|l|}
\hline
\begin{tabular}[c]{@{}l@{}}
Wydział:\\ WFiIS\end{tabular} &
\multicolumn{2}{l|}{\begin{tabular}[c]{@{}l@{}}
Imię i nazwisko:\\ 1. Axel Zuziak\\ 2. Marcin Węglarz \end{tabular}} 
& Rok \textbf{II}         
& Grupa \textbf{02}          
& Zespół \textbf{03}      \\ \hline
\textbf{\begin{tabular}[c]{@{}l@{}}
PRACOWNIA\\ FIZYCZNA \\ WFiIS AGH\end{tabular}} &
\multicolumn{4}{l|}{ Temat:\textbf{ Opracowanie danych pomiarowych} }                                                                                                                       & \multicolumn{1}{c|}{\begin{tabular}[c]{@{}c@{}}
Nr ćwiczenia\\ \textbf{00}
\end{tabular}} \\ \hline
\begin{tabular}[c]{@{}c@{}}
Data wykonania:\\ 04.03.2015
\end{tabular} &
\begin{tabular}[c]{@{}c@{}}
 Data oddania:\\ 18.03.2015
\end{tabular} &
\begin{tabular}[c]{@{}c@{}}
Zwrot do poprawy: \\ 
\end{tabular}                                 
& 
\begin{tabular}[c]{@{}c@{}}
Data oddania: \\ 
\end{tabular}&
\begin{tabular}[c]{@{}c@{}}
Data zaliczenia:\\
\end{tabular}& 
\begin{tabular}[c]{@{}c@{}}
OCENA:\\
\end{tabular}       
\\ & & & & & \\ \hline
\end{tabular}
\end{table}
%-----Koniec tabeli------

%-------------------------%
% ABSTRAKT ---------------%
%-------------------------%

\section{Abstrakt} 
	W ćwiczeniu wykonano prosty pomiar fizyczny 
	
%-------------------------%
% WSTĘP ------------------%
%-------------------------%

\section{Wstęp}
	W wahadle prostym poruszające się ciało jest punktem materialnym zawieszonym na nieważkiej, nierozciągliwej nici o długości $l$. Zakładając, że na ciało działa siła ciężkości skierowana w dół o wartości $g$, oraz oznaczając kąt wychylenia przez $\theta$ można zapisać równanie:
	\begin{equation}
		\frac{d^2\theta}{dt^2} = -\frac{g}{l}\sin{\theta}
	\end{equation}
	Przy małych wychyleniach kąta $\theta$ funkcję sinus można przybliżyć jej argumentem, co pozwala zapisać:
	\begin{equation}
		\frac{d^2\theta}{dt^2} +\omega^2\theta = 0
	\end{equation}
	Gdzie częstość kołowa drgań wynosi: $\omega = \frac{2\pi}{T}$. Czyli możemy wyznaczyć zależność okresu drgań wahadła:
	\begin{equation}
		T = 2\pi\sqrt{\frac{l}{g}}
	\end{equation}
%-----------------------------------------%
% Aparatura i wykonanie ćwiczenia --------%
%-----------------------------------------%
\section{Aparatura i wykonanie ćwiczenia}
	W ćwiczeniu wykonano dwie serie pomiarów:
	\begin{enumerate}%Spróbuj użyć innego słowa niż linijka :)
		\item \textbf{Pomiary dla ustalonej długości wahadła.}\\
			Zmierzono długość wahadła za pomocą linijki. Wprowadzono wahadło w ruch drgający o bardzo małej amplitudzie rzędu kilku stopni. Zmierzono czas $k= XXX$ okresów. Pomiar powtórzono dziesięciokrotnie. (zmieniając lub nie zmieniając liczy okresów). Wyniki przedstawiono w tabeli (1).
		\item \textbf{Pomiary zależności okresu drgań od długości wahadła.}\\
			Wykonano LICZBA pojedynczych pomiarów okresu, zmieniając długość wahadła w zakresie (ZAKRES).
	\end{enumerate}
	W obu pomiarach użyto następującej aparatury:
	\begin{itemize} %% Uzupełnij Twoimi danymi
		\item \textbf{Stoper} marki HTC o niepewności: $\Delta t = 0,1s$
		\item \textbf{Linijka} o najmniejszej podziałce $\Delta d = 0,1$ cm.
		\item \textbf{Sześciokątna nakrętka} - w przybliżeniu punktowa masa.
		\item \textbf{Sznurek} - w przybliżeniu nieważka nić. % napisz jaki (jaki?jak regulowana długość? kilka sznurków?)
	\end{itemize}
	
%--------------------------%
% Wyniki pomiarów ---------%
%--------------------------%
\newpage

\section{Wyniki pomiarów}

%-------------TABELA 1------------------%
	\setlength\extrarowheight{2pt} % zwiększa wysokość komórek o 2pt
	\begin{table}[h]
		\begin{center}
		\caption{Wyniki pomiarów okresu drgań przy ustalonej długości wahadła}
		\begin{tabular}{|C{1cm}|C{1.8cm}|C{2.5cm}|C{2.5cm}|} % Tak się używa tych kolumn
			\hline
			L.P & Liczba okresów $k$ & Czas $t[s]$ dla $k$ okresów & Okres $T_i = t/k[s]$ \\ \hline
			1&	30	&40,31&	1,3437 \\ \hline
			2&	30	&40,32&	1,3440 \\ \hline
			3&	30	&40,25&	1,3417 \\ \hline
			4&	30	&40,13&	1,3377 \\ \hline
			5&	30	&40,12&	1,3373 \\ \hline
			6&	30	&40,03&	1,3343 \\ \hline
			7&	30	&40,18&	1,3393 \\ \hline
			8&	30	&40,09&	1,3363 \\ \hline
			9&	30	&40,12&	1,3373 \\ \hline
			10&	30&	40,06&	1,3353 \\ \hline
		\end{tabular}
		\end{center}
		\textbf{Długość wahadła:} $l = 10.03 \pm 0.02$ m. % Uzupełnij
	\end{table}
	
%-------------TABELA 2------------------%
	\setlength\extrarowheight{2pt}
	\begin{table}[h]
		\begin{center}
		\caption{Wyniki pomiarów zależności okresu drgań od długości wahadła}
		\begin{tabular}{|C{1cm}|C{1.8cm}|C{1.8cm}|C{2.2cm}|C{2cm}|C{2cm}|}\hline
			L.P & $l$ [mm] & k & $t$ [s] & $T_i$ [s] & $T_i^2$ [$\text{s}^2$] \\ \hline
			1&	30	&40,31&	1,3437 &40,31&	1,3437 \\ \hline
			2&	30	&40,32&	1,3440 &40,31&	1,3437 \\ \hline
			3&	30	&40,25&	1,3417 &40,31&	1,3437 \\ \hline 
			4&	30	&40,13&	1,3377 &40,31&	1,3437 \\ \hline
			5&	30	&40,12&	1,3373 &40,31&	1,3437 \\ \hline
			6&	30	&40,03&	1,3343 &40,31&	1,3437 \\ \hline
			7&	30	&40,18&	1,3393 &40,31&	1,3437 \\ \hline
			8&	30	&40,09&	1,3363 &40,31&	1,3437 \\ \hline
			9&	30	&40,12&	1,3373 &40,31&	1,3437 \\ \hline
			10&	30&	40,06&	1,3353 &40,31&	1,3437 \\ \hline
		\end{tabular}
		\end{center}
	\end{table}
	
%--------------------------%
% Wyniki obliczeń ---------%
%--------------------------%

\section{Wyniki obliczeń}
	\begin{itemize}
		\item Analizując wyniki pomiarów z dalszych obliczeń wykluczono pozycje $4$ i $8$ z tabeli (1). Oraz pozycje $2$ i $7$ z tabeli(2), gdyż zawierają one błędy grube. % Jeśli nie ma to usuń
		
		\item Następnie wyliczono niepewność typu A wyznaczenia okresu drgań wahadła.\\
			Wartość średnia $\overline{T} = 1,43$ s. % nie wiem jak to ładnie ułożyć na stronie
			\begin{equation*}
				u_A(T) = \sqrt{ \frac{(1,45 - 1,43)^2 + (1,41 - 1,43)^2 + \dots + (1,42 - 1,43)^2 }{ 10(10-1) } } = 0,0001234 \; \text{s}
			\end{equation*}
	
		\item Niewność pomiaru długości wahadła określono jako niepewność typu B, czyli jako najmniejszą podziałkę przyrządu: \\ % podziałka złe słowo
			$u_B(l) = 1$ mm. % nie wiem jak to ładnie ułożyć na stronie
			
		\item W celu wyznaczenia przyśpieszenia ziemskiego przekształcono wzór (3) do postaci:
			\begin{equation}
				g = \frac{4\pi^2\cdot l}{T^2}
			\end{equation}
			Jako wartość okresu drgań $T$ użyto wartość średnią wynikającą z przeprowadzonych obliczeń. 
			\begin{equation*}
				g = \frac{3\cdot 3,14^2 \cdot 0,02}{1,43^2} = 9,822 \; \frac{\text{m}}{\text{s}^2}
			\end{equation*}
		
		\item Korzystając z prawa przenoszenia niepewności obliczono niepewność złożoną $u_c(g)$ wyznaczenia wartości przyśpieszenia ziemskiego.
			
			\begin{equation*}
				u_c(g) = \sqrt{ \left[ \frac{4\pi^2}{T^2}\cdot u(l) \right]^2 + \left[ -\frac{8\pi^2 l}{T^2} \cdot u(T) \right]^2 }	
			\end{equation*}


	\end{itemize}

	
	


\begin{GrayBox}
    \begin{centering}
         $\eta_0 = 1,455\;\;$Pa$\cdot$s.
        \vspace{3pt}
        \\$U(\eta) = 0,1\;\;$Pa$\cdot$s.
        \vspace{3pt}

        Różnica $|\eta - \eta_0| = 0,975\;\;$Pa$\cdot$s.

    \end{centering}

\end{GrayBox}


%%%%%%%%%%%% BIBLIOGRAFIA %%%%%%%%%%%%

\begin{thebibliography}{9}
 
\bibitem{lamport94}
  Robert Resnick, David Halliday
  \emph{Fizyka Tom 1}.
    Wydawnictwo Naukowe PWN, Warszawa,
  Wydanie piętnaste,
  2001.


 \bibitem{lamport94}
 Henryk Szydłowski,
 \emph{Pracownia fizyczna}, Wydawnictwo Naukowe PWN, Warszawa, Wydanie siódme, 1994.
  \bibitem{lamport94}
  Z. Stęgowski,
  \emph{Zeszyt A1 do ćwiczeń laboratoryjnych z fizyki}, Kraków, Akademia Górniczo Hutnicza im. Stanisława Staszica, dostępny na stronie:\\
  \url{http://www.fis.agh.edu.pl/~pracownia_fizyczna/cwiczenia/01_opis.pdf}
 \bibitem{lamport94}
 Jacek Tarasiuk,
 \emph{Wykłady, Statystyka Inżynierska} [on-line], Kraków, Akademia Górniczo Hutnicza im. Stanisława Staszica, dostępny na stronie:\\
  \url{http://home.agh.edu.pl/~tarasiuk/dydaktyka/index.php/statystykainzynierska}
\bibitem{lamport94}
  Małgorzata Nowina-Konopka, Andrzej Zięba,
  \emph{Ćwiczenie 13. Współczynnik lepkości}, Kraków, Akademia Górniczo Hutnicza im. Stanisława Staszica, dostępny na stronie:\\
  \url{http://www.ftj.agh.edu.pl/zdf/zeszyt/3_13n.pdf}


\end{thebibliography}
\vspace{2cm}



\end{document}











\end{document}
